\documentclass[11pt]{article}

\usepackage[margin=1.05in]{geometry}
\usepackage{lmodern}
\usepackage{amsmath,amssymb}
\usepackage{booktabs}
\usepackage{enumitem}
\usepackage{xcolor}
\usepackage[expansion=false]{microtype}
\usepackage[colorlinks=true,linkcolor=blue!55!black,urlcolor=blue!55!black]{hyperref}

\newcommand{\code}[1]{\texttt{#1}}
\DeclareMathOperator*{\argmax}{arg\,max}
\newcommand{\Rset}{\mathcal{R}}

\title{\bfseries A Two-Step Model for Startup Funding:\\[2pt]
\large Sector-level self-exciting counts and within-sector ranking}
\author{Alberto Gennaro}
\date{24 June 2026}

\begin{document}
\maketitle

\section{Overview}

The task is to predict, from weekly funding data, which startup is funded next. The
model factorises this into two layers that match the structure of the data:

\begin{enumerate}[leftmargin=1.6em,itemsep=2pt]
\item a \emph{sector} layer that predicts \emph{where and when} funding events occur ---
a weekly multivariate self-exciting count model over $M$ sectors, with positive lagged
self- and cross-excitation and market covariates;
\item a \emph{within-sector} layer that predicts \emph{which} startup receives a given
sector event --- a conditional ranker over the live startups in that sector, with a
recency (``cooldown'') term that suppresses startups that funded recently.
\end{enumerate}

The decomposition is exact under a marked-process view: the probability that the next
event is a particular startup is the sector intensity times the within-sector choice
probability. It also separates a low-dimensional, identifiable counting problem (the
sector layer) from a high-cardinality, dynamic-risk-set problem (the ranker). Both
layers are fit by maximum likelihood. The implementation is in
\code{hawkes\_calibration/sector\_ranker.py}; the end-to-end synthetic backtest is in
\code{sector\_backtest.py} and \code{experiments/exp14\_sector\_ranker.py}. Numerical
results quoted in Section~\ref{sec:study} are from the project's learning report
(five-seed synthetic backtest); no experiment is re-run here.

\section{Data and notation}

Time is discrete (weeks) $t=1,\dots,T$. There are $M$ sectors and a roster of startups
that changes over time. Observed quantities:
\begin{itemize}[leftmargin=1.4em,itemsep=1pt]
\item $Y_{s,t}\in\mathbb{N}$: number of funding events in sector $s$ in week $t$.
\item $X_t\in\mathbb{R}^{p}$: market-level covariates (e.g.\ a risk-on regime).
\item $Z_{j,t}\in\mathbb{R}^{q}$: features of startup $j$ at week $t$ (quality,
founder/investor signal, traction, a sector-momentum feature).
\item $\Rset_{s,t}$: the live risk set --- the startups currently in sector $s$ and
active at week $t$. Its size changes as startups enter or exit.
\item events $(t,s,i^\ast)$: in week $t$, sector $s$, startup $i^\ast$ was funded.
\end{itemize}

\section{Data: sources, features, and run modes}
\label{sec:data}

\paragraph{Sources.} Two licensed (non-public) providers supply the inputs.
\emph{FactSet Standardized Economics} provides the macro covariates --- interest
rates, inflation and unemployment --- via the series \code{FFUNDT}, \code{FFUND} (US
Fed Funds target / effective), \code{CPI\_RAW}, \code{CPI\_YOY} (US CPI level /
year-over-year) and \code{RUNEMP} (unemployment rate). \emph{PitchBook} provides the
event stream and company attributes through its Deals, Company and Investors tables.
Both are subscription products and are not publicly available. The macro series,
however, have free public equivalents on FRED (\code{FEDFUNDS} or \code{EFFR},
\code{CPIAUCSL}, \code{UNRATE});\footnote{\url{https://fred.stlouisfed.org}} for
deal-level PE/VC activity the closest free substitutes are the Crunchbase free tier and
SEC Form~D filings (public, with a short lag). Because the deal data is proprietary, the
pipeline runs on synthetic data by default and reads user-supplied CSVs in a real mode.

\paragraph{From raw fields to model inputs.} A deal's \code{primary\_industry\_sector}
defines its sector $s$ and its \code{deal\_date} its week $t$; the sector count
$Y_{s,t}$ in \eqref{eq:sector} is the number of deals in $(s,t)$.
\begin{itemize}[leftmargin=1.4em,itemsep=2pt]
\item \textbf{Sector covariates $X_{s,t}$} (the baseline of \eqref{eq:sector}): the
standardized macro series and their first differences (shared across sectors); plus
per-sector weekly aggregates from the Deals table, lagged one week and standardized ---
total amount raised ($\sum$\code{deal\_size}), amount per event (total\,/\,count), mean
\code{deal\_size}, the up/down/flat round ratio (\code{vc\_round\_up\_down\_flat}), mean
\code{pre\_money\_valuation} and \code{post\_valuation}, mean \code{revenue} and
\code{ebitda}, and mean \code{number\_of\_employees}. These are known before week $t$,
so they are valid regressors; the self- and cross-excitation themselves enter through
the lagged counts in \eqref{eq:sector}.
\item \textbf{Startup features $Z_{j,t}$} (the score \eqref{eq:score}): company- and
deal-level attributes --- age (from \code{year\_founded}), \code{total\_raised\_to\_date},
\code{last\_financing\_size}, round/stage (\code{vc\_round}), \code{number\_of\_employees},
\code{revenue}, \code{ebitda}, \code{revenue\_growth\_since\_last\_deal},
\code{total\_investor\_ownership}, \code{business\_status} and \code{hq\_global\_region};
and dynamic features --- weeks since the last financing (the recency / cooldown
$c^{(K)}_{j,t}$), the number of prior rounds, and a recent-sector-activity ``momentum''
feature. Continuous features are standardized; categoricals are one-hot encoded.
\item \textbf{Risk set $\Rset_{s,t}$}: a company enters at its first appearance (or
\code{year\_founded}) and leaves when \code{current\_financing\_status} or
\code{ownership\_status} indicates an acquisition, IPO or closure; the live set in sector
$s$ at week $t$ is the candidate pool the ranker normalises over.
\end{itemize}

\paragraph{Run modes.} The estimation of Sections~\ref{sec:step1}--\ref{sec:step2} is
identical in both modes; only the source of the raw columns differs.
\begin{itemize}[leftmargin=1.4em,itemsep=2pt]
\item \emph{Synthetic} (default). The generator
(\code{simulate\_synthetic\_startup\_market}) populates the raw columns above with
plausible distributions --- log-normal, positively correlated deal sizes and valuations;
categorical rounds; AR(1) macro series; random founding dates and entries --- and the
same aggregation/feature pipeline builds $X$, $Z$, $Y$ and the risk sets. No proprietary
data is required.
\item \emph{Real} (\code{-{}-mode real -{}-data-dir DIR}). Two CSV files are read from
\code{DIR}: \code{factset\_macro.csv} (a date column plus the macro series
\code{FFUNDT}, \code{FFUND}, \code{CPI\_RAW}, \code{CPI\_YOY}, \code{RUNEMP}, in long or
wide form) and \code{pitchbook\_deals.csv} (one row per deal, with at least
\code{company\_id}, \code{deal\_id}, \code{deal\_date}, \code{primary\_industry\_sector},
\code{deal\_size}, \code{vc\_round}, \code{vc\_round\_up\_down\_flat},
\code{pre\_money\_valuation}, \code{post\_valuation}, \code{total\_raised\_to\_date},
\code{number\_of\_employees}, \code{revenue}, \code{ebitda}, \code{year\_founded},
\code{current\_financing\_status}, \code{hq\_global\_region}). Optional
\code{pitchbook\_companies.csv} and \code{pitchbook\_investors.csv}, joined on
\code{company\_id} and \code{investor\_id}, add further company and investor features.
Preprocessing maps \code{deal\_date} to a week index, \code{primary\_industry\_sector}
to $s$, forward-fills the macro series onto the weekly grid, aggregates deals into the
$(s,t)$ counts and sector aggregates, and constructs the startup features and dynamic
risk sets.
\end{itemize}

\paragraph{Feature reduction.} The raw PitchBook schema is wide --- dozens of deal- and
company-level fields, many sparse, skewed or collinear --- which inflates the variance of
the ranker's sector-specific weights and invites over-fitting. Preprocessing therefore
(i)~log-transforms the monetary fields (\code{deal\_size}, valuations, \code{revenue}),
standardises continuous features, groups rare categories and adds explicit missingness
indicators; and (ii)~reduces dimension before fitting --- a truncated SVD / PCA on the
standardised continuous block keeps the top $r$ components (those explaining, say, $90\%$
of variance), and the sector-specific ranker weights are given a low-rank form
$u_s = U_s V$ with $r\ll q$, cutting the parameter count from $Mq$ to $(M+q)r$. This is
the low-rank / embedding device used to scale relational point-process models
(Know-Evolve, DyRep) and keeps the high-cardinality feature set from overwhelming the
conditional-logit fit.

\section{Step 1: the sector-level self-exciting count model}
\label{sec:step1}

\paragraph{Model.} Sector counts follow a Poisson log-linear autoregression --- a
discrete-time multivariate Hawkes process,
\begin{equation}
Y_{s,t}\sim\mathrm{Poisson}(\Lambda_{s,t}),\qquad
\log\Lambda_{s,t}=a_s+\beta_s^\top X_t+\sum_{r=1}^{M}\sum_{\ell=1}^{L}b_{s,r,\ell}\,Y_{r,t-\ell},
\label{eq:sector}
\end{equation}
with intercepts $a_s$, covariate loadings $\beta_s$, and nonnegative lag coefficients
$b_{s,r,\ell}\ge 0$. The diagonal terms $b_{s,s,\ell}$ are self-excitation (a sector that
funded recently funds more); the off-diagonal terms $b_{s,r,\ell}$ ($r\neq s$) are
cross-sector contagion. An optional adjacency mask forces chosen entries to zero. The
nonnegativity constraint keeps the rate positive and the excitation interpretable.

\paragraph{Estimation.} The parameters $\theta=(a,\beta,b)$ are estimated by penalised
Poisson maximum likelihood. Up to a constant the objective is
\begin{equation}
\min_{\theta:\,b\ge 0}\ \sum_{t=L+1}^{T_{\mathrm{tr}}}\sum_{s=1}^{M}
\big(\Lambda_{s,t}-Y_{s,t}\log\Lambda_{s,t}\big)
+\tfrac{\lambda}{2}\big(\lVert\beta\rVert^2+\lVert b\rVert^2\big),
\label{eq:sectorfit}
\end{equation}
solved with a box-constrained quasi-Newton method (L-BFGS-B); the gradient is available
in closed form, and the intercepts are left unpenalised. The problem is convex, so the
fit is unique. Held-out scoring uses the one-step rate
$\Lambda_{s,t}$ formed from the realised lagged counts.

\paragraph{Stability.} Equation~\eqref{eq:sector} is stable (the mean count stays
finite) only if the excitation has spectral radius below one. The estimator constrains
$b\ge 0$ but not the spectral radius, so a fitted model can be supercritical; in the
synthetic study a fit reached spectral radius $\approx 3.7$. This does not affect
one-step held-out scoring, but it makes \emph{forward simulation} explode and must be
corrected (project the fitted excitation to spectral radius below one, and cap the inner
event loop) before simulating paths.

\section{Step 2: the within-sector survival model with an outside option}
\label{sec:step2}

Given a sector event from Step~1, we model which startup receives it as a discrete-time
competing-risks survival problem over the live risk set, with one addition that matches a
positives-only deals feed: an explicit \emph{outside} alternative --- ``a firm not in our
dataset'' --- so an event that goes to a company we do not track is not forced onto an
observed candidate.

\paragraph{Candidates and the outside option.} Each tracked candidate $j\in\Rset_{s,t}$
(the live startups in sector $s$ that we follow) has a log-hazard
\begin{equation}
q_{j,t}=(w_0+u_s)^\top Z_{j,t}+\eta_s\,c^{(K)}_{j,t},
\label{eq:score}
\end{equation}
with global feature weights $w_0$, sector deviations $u_s$, and a cooldown term
($c^{(K)}_{j,t}$ is startup $j$'s funding count over the past $K$ weeks; $\eta_s\le0$
makes a recently funded startup less likely to fund again --- self-inhibition). The
outside alternative $0$ carries a log-hazard built from \emph{sector-level} features only,
since it has no company covariates,
\begin{equation}
q_{0,s,t}=b_0+\gamma^\top g_{s,t},\qquad
g_{s,t}=\big(\log(1+|\Rset_{s,t}|),\ \text{recent sector activity}\big).
\label{eq:outside}
\end{equation}
The recipient is drawn from $\Rset_{s,t}\cup\{0\}$ by a softmax over the augmented set,
\begin{equation}
\Pr(i\mid s,t)=\frac{\exp(q_{i,t})}{\exp(q_{0,s,t})+\sum_{j\in\Rset_{s,t}}\exp(q_{j,t})},
\label{eq:softmax}
\end{equation}
$i$ ranging over the tracked candidates and the outside option. Dropping $q_0$ recovers
the plain risk-set ranker.

\paragraph{Why the outside option fits the data.} A deals database records only firms that
raised; the universe of firms that could have raised is unobserved, and company features
are typically available only at a firm's last round. The outside option turns this from a
defect into a modelled quantity. At training, events to firms outside the tracked universe
(a watch-list / portfolio) are labelled ``outside'' and supply the signal for
$(b_0,\gamma)$; at prediction, $\Pr(\text{outside}\mid s,t)$ is the calibrated probability
that the next funded firm is one we do not yet track --- a useful output in itself, since a
high value says ``widen the watch-list.'' Firms never seen carry no features and are
simply not candidates; they fall under the outside mass.

\paragraph{Estimation.} The parameters $(w_0,u,\eta,b_0,\gamma)$ are estimated by
conditional maximum (partial) likelihood,
\begin{equation}
\min_{\eta\le0}\ \sum_{(t,s,i^\ast)}\Big(\log\big[\exp(q_{0,s,t})+\!\!\sum_{j\in\Rset_{s,t}}\!\!\exp(q_{j,t})\big]-q_{i^\ast,t}\Big)+\text{(ridge penalties)},
\label{eq:rankfit}
\end{equation}
with $i^\ast$ the funded candidate or the outside option. The objective is convex and
box-constrained (cooldown nonpositive); its gradient is the softmax residual. The
estimator (\code{fit\_startup\_survival}) is numpy-only, with SciPy used when available.

\paragraph{Relation to standard models.} Equation~\eqref{eq:rankfit} is the log-sum-exp
partial likelihood of a discrete choice over the risk set --- McFadden's conditional logit
\emph{with an outside good}, equivalently the Cox proportional-hazards partial likelihood
(competing risks: each candidate plus an outside cause) and the Plackett--Luce choice
model. The cooldown is a time-varying covariate; the negative-constrained $\eta_s$ and the
self-exciting sector layer below it are what separate this model from the feature-only
benchmark of Section~\ref{sec:bench}. On the synthetic market the analytic gradient
matches finite differences to $10^{-8}$, the fitted cooldown is nonpositive, the outside
option discriminates off-list events (area under ROC $\approx0.68$), and the
within-universe ranking beats chance (top-5 $\approx0.70$ over $\sim$7 candidates); see
\code{tests/test\_sector\_survival.py}.

\section{Producing a ranked watch-list}

Within a sector, startups are ranked by \eqref{eq:softmax} (or by the score
\eqref{eq:score}). For an overall ``who is funded next'' list across sectors, the
marked-process factorisation gives a propensity
\begin{equation}
\text{score}_{j,t}\ \propto\ \Lambda_{s(j),t}\cdot\Pr(j\mid s(j),t),
\label{eq:marked}
\end{equation}
the sector intensity from Step~1 times the within-sector choice probability from Step~2.
Ranking quality is measured on held-out events by the top-$k$ hit rate (is the funded
startup in the top $k$), the mean reciprocal rank (MRR), and the held-out negative
log-likelihood, each compared against a random-over-the-risk-set baseline and against the
\emph{oracle} ranker that uses the true parameters (the best achievable given the
stochastic draw).

\section{Study case: a synthetic market}
\label{sec:study}

\paragraph{Data-generating process.} This is the synthetic run mode of
Section~\ref{sec:data}; the real mode applies the identical estimation and evaluation to
user-provided FactSet/PitchBook CSVs. The simulator
(\code{simulate\_synthetic\_startup\_market}) mirrors the two-layer model. Sector counts
are drawn from \eqref{eq:sector} with a sparse, strongly diagonal positive excitation
matrix (self-excitation per sector plus two cross-sector links) and AR(1) market
covariates. Each sector event is assigned to a live startup through the risk-set softmax
\eqref{eq:softmax} with a \emph{negative} true cooldown, so the ground truth has both
self-excitation (sectors) and self-inhibition (startups). About a quarter of startups
enter after the start, so risk-set sizes change over time; startup features include a
quality signal and a sector-momentum-like feature.

\paragraph{Protocol.} A market of $M=11$ sectors and roughly $35$ startups per sector is
simulated for $T=180$ weeks; the first $120$ weeks train both layers, the last $60$ are
held out. Reported figures are averaged over five seeds.

\paragraph{Results} (from the project learning report). The sector layer improves the
held-out one-step Poisson log-likelihood over a historical-mean baseline. The ranker, on
held-out weeks with about $35$ candidates per pick, attains

\begin{table}[h]
\centering\small
\begin{tabular}{@{}l c c c@{}}
\toprule
& risk-set ranker & oracle (true parameters) & random \\
\midrule
top-5 hit rate & $\approx 0.49$ & $0.54$ & $0.14$ \\
MRR & \multicolumn{1}{c}{$\approx 0.88$ of oracle} & --- & --- \\
\bottomrule
\end{tabular}
\end{table}

The ranker is three to six times better than chance and recovers roughly $80$--$92\%$ of
the oracle across seeds. The remaining gap to the oracle is the irreducible stochasticity
of the funded-startup draw over the risk set, not model error. The sector-layer
spectral-radius caveat of Section~\ref{sec:step1} applies to forward simulation, not to
this one-step held-out scoring.

\section{Benchmark: a feature-based ranker without explicit self-excitation}
\label{sec:bench}

The two-step model encodes self-excitation and self-inhibition as \emph{mechanisms}: the
nonnegative lagged excitation matrix in \eqref{eq:sector}, and the sign-constrained
cooldown coefficient in \eqref{eq:score}. A natural benchmark removes these mechanisms
and instead supplies engineered \emph{features} that proxy for the same effects, to test
whether the explicit structure earns its place.

\paragraph{Construction.} Keep the risk-set evaluation identical (same held-out events,
same live risk sets, same metrics), but:
\begin{itemize}[leftmargin=1.4em,itemsep=2pt]
\item \textbf{No sector self-excitation.} Replace Step~1 by a feature-only rate --- a
Poisson GLM on the covariates $X_t$ and an exogenous trailing-activity feature (e.g.\ a
moving average of recent sector counts treated as an input, not an estimated kernel), or
simply the historical-mean rate. There is no estimated lagged excitation matrix.
\item \textbf{No explicit self-inhibition.} Rank startups with a survival / discrete-time
hazard model (a Cox proportional-hazards or DeepHit-style logistic-hazard ranker) in which
``weeks since last funding'' and a ``recent sector momentum'' enter as ordinary
covariates with free (unconstrained) coefficients. The recency feature is meant to absorb
the cooldown; the momentum feature to absorb the sector self-excitation.
\end{itemize}
This benchmark is exactly a standard survival ranker with hand-built momentum and recency
features. The package already provides the survival machinery for it
(\code{SURVIVAL\_RANKING.md}, \code{experiments/exp16\_survival.py}: Kaplan--Meier, a
discrete-time hazard with a ranking loss, concordance and Recall@$k$).

\paragraph{What the comparison isolates.} Run the two-step model and the feature
benchmark on the same backtest and compare top-$k$, MRR and the time-dependent
concordance index. Two readings are possible and both are informative:
\begin{itemize}[leftmargin=1.4em,itemsep=2pt]
\item If the explicit two-step model wins, the \emph{mechanism} (a self-exciting sector
kernel; a sign-constrained, risk-set-normalised cooldown) carries information that
hand-built features do not --- chiefly the cross-sector contagion in \eqref{eq:sector} and
the dynamic risk-set normalisation in \eqref{eq:softmax}, neither of which a per-startup
feature reproduces.
\item If the benchmark matches it, the effects are well summarised by recency and
momentum features, and the simpler, cheaper survival ranker suffices.
\end{itemize}
The honest prior is that the gap on the \emph{who} question (Step~2) is modest, because
the cooldown is itself a recency feature; the larger differentiator is the \emph{where and
when} question (Step~1), where cross-sector excitation and the contagion dynamics are hard
to capture with a single per-sector momentum feature.

\section{Modeling assumptions and alternatives}
\label{sec:alternatives}

\paragraph{What the data supports.} The two layers make very different demands on the
data. The sector layer needs only weekly counts per sector and macro covariates, both
well observed; it is robust. The within-sector ranker is more demanding: it assumes that
at each week the live risk set of candidate startups can be enumerated and that every
candidate has features. Two facts about PE/VC data strain this. First, company covariates
(revenue, valuation, employees) are typically refreshed only \emph{at a financing event},
so between rounds they are stale and for never-funded firms they are absent --- features
arrive on an event-driven, irregular schedule, not continuously. Second, a deals database
records \emph{positives} (firms that raised); the universe of firms that could have raised
but did not --- the negatives that define the risk set --- is largely unobserved.

\paragraph{When the model is appropriate.} The conditional risk-set ranker is sound when
``candidates'' is a defined, feature-complete universe that we track --- a watch-list, a
fund's screened pipeline, or a portfolio --- rather than the entire private market. There
the risk set is observed, features are maintained, and the softmax over candidates is well
grounded; that is the recommended deployment.

\paragraph{Alternatives.} If only the positives-only deals table is available, or features
are sparse, three reformulations fit the data better.
\begin{itemize}[leftmargin=1.4em,itemsep=2pt]
\item \textbf{The outside-option survival model (adopted as Step~2).} The model of
Section~\ref{sec:step2} is exactly this reformulation: a competing-risks partial
likelihood over the tracked risk set with an explicit outside cause, so it never has to
enumerate the unobserved negatives and it reports $\Pr(\text{outside})$ directly. The
natural extension is a fuller per-firm \emph{time-to-event} version (Cox / DeepHit over
weeks since the last round, with right-censoring and last-observation-carried-forward
features), with the sector intensity multiplying the per-firm hazard,
$\Lambda_{s,t}\cdot\mathrm{hazard}_{j,t}$.
\item \textbf{Lean on always-observable signals.} The strongest, always-computable
features come from the event stream itself --- recency (weeks since last round), round
number, sector momentum --- consistent with the empirical finding that financing recency
dominates next-round prediction. A model built primarily on these, with hard covariates as
optional refinements, is robust to the missingness above.
\item \textbf{Use investor activity.} The Investors table is observed and informative:
``which firm is funded next'' is partly ``which firm an active investor selects.'' A
bipartite investor--company dynamic-graph model (Know-Evolve / DyRep style) ranks companies
by predicted investor--company edge formation and avoids relying on company financials.
\end{itemize}
In short, the sector layer is well-supported and the two-step architecture is sound, but
the within-sector ranker should either be scoped to a tracked candidate universe or recast
as a per-firm survival model; the cold-start of never-seen firms and the unobserved
negatives are the binding constraints, not the choice of estimator.

\section{Limitations}

\begin{itemize}[leftmargin=1.3em,itemsep=2pt]
\item Results are on synthetic data drawn from the same two-layer family the model
assumes; the real test is a back-test on observed weekly funding data.
\item The sector excitation is not spectral-radius-constrained during fitting; forward
simulation requires projecting to the stable region.
\item The decay structure (number of lags $L$, cooldown window $K$) is fixed, not
selected; the cross-excitation entries are weakly identified when sector counts are low.
\item The benchmark of Section~\ref{sec:bench} is specified but not yet run; it is the
recommended next experiment.
\end{itemize}

\section{References}

\begin{enumerate}[leftmargin=1.6em,itemsep=2pt]
\item Rizoiu, Soen, Li, Calderon, Dong, Menon, Xie (2022). Interval-censored Hawkes
processes. JMLR 23(1). \url{https://www.jmlr.org/papers/v23/21-0917.html}
\item Hawkes (1971). Spectra of some self-exciting and mutually exciting point processes.
Biometrika 58(1).
\item Fokianos, Rahbek, Tj\o stheim (2009). Poisson autoregression. JASA 104(488):1430--1439.
(discrete-time self-exciting counts.)
\item McFadden (1974). Conditional logit analysis of qualitative choice behavior. In
\emph{Frontiers in Econometrics}. (the risk-set choice model.)
\item Cox (1972). Regression models and life-tables. JRSS-B 34:187--220. (partial
likelihood over the risk set.)
\item Plackett (1975). The analysis of permutations. Applied Statistics 24(2):193--202.
(ranking / choice likelihood.)
\item Kaplan \& Meier (1958). Nonparametric estimation from incomplete observations. JASA 53.
\item Lee, Zame, Yoon, van der Schaar (2018). DeepHit. AAAI.
\url{https://cdn.aaai.org/ojs/11842/11842-13-15370-1-2-20201228.pdf}
\item Antolini, Boracchi, Biganzoli (2005). A time-dependent discrimination index for
survival data. Statistics in Medicine 24(24).
\end{enumerate}

\noindent\emph{Related documents:} \code{SURVIVAL\_RANKING.md} (the survival ranker in
full); \code{NEXT\_FUNDING\_DESIGN.md} (the application design);
\code{paper/investment\_case\_study.pdf} (the continuous-time Hawkes/MBPP foundation).

\end{document}
